% ===================== page 1 (title page) =====================
\printedpage{--}
% no page number printed on this page
% head: (none)

\begin{center}
{\LARGE\sffamily THE DERIVATORS}

\bigskip

\textbf{ALEXANDRE GROTHENDIECK}

\bigskip\bigskip

\textbf{Chapter III}

\medskip

\textbf{External Hom in derivators}
\end{center}

\clearpage

% ===================== page 2 (editorial note) =====================
\printedpage{--}
% no page number printed on this page
% head: (none)

This text was deciphered and transcribed in \LaTeXe{} by M. Künzer. It was edited by M. Künzer, J. Malgoire, and G. Maltsiniotis. The transcription is as faithful as possible to the manuscript. For the few obvious corrections, or rare comments by the editors, as well as for the original numbering of the manuscript pages, typewriter characters \texttt{[typewriter]} in brackets are used. A question mark in brackets means that one is not sure of the preceding word.

This edition is provisional. Remarks, comments and corrections are welcome. Send a message to:

\textbf{[e-mail address removed]}
\begin{flushright}
G. Maltsiniotis\qquad\qquad
\end{flushright}

\clearpage

% ===================== page 3 (printed 1) =====================
\printedpage{1}
% head: A. Grothendieck : Dérivateurs, Ch. III, Édition des Universités de Montpellier II et Paris VII
% foot: Version 2 février 2002

\msp{page 1}

\section*{1\quad Relation of derivators with the category $\Hot$}

Let $\D$ be a derivator, $\mathcal{A}$ its fundamental category,
\[
\mathcal{A} = \D(e).
\]
I propose to define, if it can be done, functors
\[
\left\{
\begin{array}{rccccccl}
\Hot & \times & \mathcal{A} & \xrightarrow{\pi_{\D}} & \mathcal{A} & (\xi, M) & \longmapsto & \xi \times M \ \text{ or } \ \xi \otimes M \quad \text{(depending on the case)}\\
\Hot^o & \times & \mathcal{A} & \xrightarrow{\varpi_{\D}} & \mathcal{A} & (\xi, M) & \longmapsto & \hm(\xi, M).
\end{array}
\right.
\tag*{(1)}
\]
These functors must be interchanged by duality: if one takes the functors $\pi^o$, $\varpi^o$ on the opposite categories, one finds
\[
\left\{
\begin{array}{rcccc}
\Hot & \times & \mathcal{A}^o & \xrightarrow{\varpi^o_{\D}} & \mathcal{A}^o\\
\Hot^o & \times & \mathcal{A}^o & \xrightarrow{\pi^o_{\D}} & \mathcal{A}^o,
\end{array}
\right.
\tag*{(1$'$)}
\]
and one must have
\[
\left\{
\begin{array}{rcl}
\varpi^o_{\D} & = & \pi_{\D^o}\\
\pi^o_{\D} & = & \varpi_{\D^o}.
\end{array}
\right.
\]
To tell the truth, considering (1) as functors
\[
\begin{array}{rcl}
\Hot & \xrightarrow{\pi_{\D}} & \uHom(\mathcal{A}, \mathcal{A})\\
\Hot^o & \xrightarrow{\varpi_{\D}} & \uHom(\mathcal{A}, \mathcal{A}),
\end{array}
\]
one expects these functors to extend to
\[
\left\{
\begin{array}{rcl}
\Hot & \xrightarrow{\pi_{\D}} & \uHom(\D, \D)\\
\Hot^o & \xrightarrow{\varpi_{\D}} & \uHom(\mathcal{A}, \mathcal{A})
\end{array}
\right.
\tag*{(2)}
\]
\texttt{[rather twice} $\uHom(\D, \D)$\texttt{]}. This means, essentially, that for a given $\xi \in \mathrm{Ob}\,\Hot$, the operations $\xi \times_{\D} M$ and $\uHom_{\D}(\xi, M)$ \texttt{[rather} $\hm_{\D}(\xi, M)$\texttt{]} are defined not only in $\mathcal{A} = \D(e)$, but in all the $\D(I)$, and this in such a way as to commute with the functors $f^*$:

\msp{page 2}
If
\[
f : I' \longrightarrow I,
\]
\printedpage{2}
% head: A. Grothendieck : Dérivateurs, Ch. III, Édition des Universités de Montpellier II et Paris VII
% foot: Version 2 février 2002
then for $F \in \mathrm{Ob}\,\D(I)$, one must have canonical isomorphisms
\[
\left\{
\begin{array}{rcl}
f^*(\xi \times_{\D} M) & \simeq & \xi \times_{\D} f^*(M)\\
f^*\hm_{\D}(\xi, M) & \simeq & \hm_{\D}(\xi, f^*(M)),
\end{array}
\right.
\tag*{(3)}
\]
with the usual compatibilities for a composite $fg$, relative to $I'' \xrightarrow{g} I' \xrightarrow{f} I$.

But I prefer to work in $\mathcal{A}$, which should imply the similar constructions in all the $\D(I)$, by applying what has been said to the derivator $\D_I$ if need be.

The fundamental algebraic intuition is this, denoting by
\[
X \longmapsto \mathrm{hot}(X) : (\mathrm{Cat}) \longrightarrow (\Hot)
\]
the localization functor from (Cat) to (Hot), obtained by inverting the arrows in (Cat) which are ``weak equivalences''. Suppose $X$ is in Diag, let
\[
p : X \longrightarrow e
\]
be the functor with values in the one-point category. We want
\[
\left\{
\begin{array}{rcccl}
\mathrm{hot}(X) \times_{\D} M & \simeq & p_!(p^*(M)) & = & H^{\D}_{\bullet}(X, M)\\
\hm_{\D}(\mathrm{hot}(X), M) & \simeq & p_*(p^*(M)) & = & H^{\bullet}_{\D}(X, M)
\end{array}
\right.
\tag*{(4)}
\]
(canonical isomorphisms).

\msp{page 3}
This would give a principle of definition for the operations sought. A difficulty comes from the fact that in general, Diag is not equal to (Cat), and that the homotopy types coming from Diag do not give (even up to isomorphism) \textit{all} homotopy types. Intuitively, one finds only ``small enough'' homotopy types. But one can localize Diag with respect to weak equivalences (it is better to say ``quasi-isomorphisms''), one then finds a category $(\Hot)_0$, and a canonical functor
\[
\Hot_0 \longrightarrow \Hot,
\]
if we are lucky, this functor will be fully faithful, and we shall try to describe the functors (1) or (2) by replacing Hot in them by the full subcategory of ``small'' homotopy types. One can also forget (temporarily!) Hot, and work with $\Hot_0$. This is what we are going to do for the moment.

In order to be able to define the functors
\[
\left\{
\begin{array}{lccccccl}
\Hot_0 & \times & \mathcal{A} & \longrightarrow & \mathcal{A} & (\xi, F) & \longmapsto & \xi \times_{\D} F\\
\Hot_0^o & \times & \mathcal{A} & \longrightarrow & \mathcal{A} & (\xi, F) & \longmapsto & \hm_{\D}(\xi, F),
\end{array}
\right.
\tag*{(5)}
\]
I note that for $F$ \texttt{[rather} $M$\texttt{]} fixed, $X$ variable in Diag, the
\printedpage{3}
% head: A. Grothendieck : Dérivateurs, Ch. III, Édition des Universités de Montpellier II et Paris VII
% foot: Version 2 février 2002
\msp{page 4}
last members in (4)
\[
H^{\D}_{\bullet}(X, M), \qquad H^{\bullet}_{\D}(X, M)
\tag*{(6)}
\]
are indeed, the one covariant, the other contravariant in $X$. Thus there is , up to unique isomorphism, a single pair of bifunctors (5) satisfying (bifunctorially) the isomorphisms (4), if and only if the expressions (6), considered as functorial in $X$, transform quasi-isomorphisms into isomorphisms. This is therefore equivalent to the validity of the % [sic] "il y a ," -> "il y a,", espace imprimée avant la virgule

\textbf{Theorem} (?) (\footnote{New axiom?}). \textit{Let $f : X' \longrightarrow X$ be in Diag. If $f$ is a quasi-isomorphism, it is a $\D$-equivalence.}

When $X$ is the one-point category, this is nothing other than axiom Der 3 (p. 67). I do not know whether it is a consequence of axioms Der 1 to Der 5 (d). If this is not the case, one would have to add it to the axioms, I think. I have not even proved that an equivalence of categories in Diag (in the sense of categories!) is a $\D$-equivalence!

Let us note the \textit{adjunction formula}
\[
\mathrm{Hom}(\xi \times_{\D} M, N) \xrightarrow{\sim} \mathrm{Hom}(M, \hm_D(\xi, N)).
\tag*{(7)}
\] % [sic] \hom_D (D italique maigre) -> \hom_{\mathbf{D}}, le dérivateur est noté en gras partout ailleurs
If $\xi$ comes from $X \in \mathrm{Diag}$, then this reads
\[
\mathrm{Hom}(p_!p^*M, N) \xrightarrow{\sim} \mathrm{Hom}(M, p_*p^*N),
\tag*{(8)}
\]
\msp{page 5}
now by adjunction both sides are canonically isomorphic to
\[
\mathrm{Hom}_{\D(X)}(p^*M, p^*N).
\]

\textbf{Example.} Taking for $\xi$ the homotopy type of the standard circle (represented e.g. by the ordered set
\[
\begin{tikzcd}[row sep=small, column sep=small]
 & b & \\
0 \arrow[ur] \arrow[dr] & & 0' \arrow[ul] \arrow[dl] \ , \\
 & a &
\end{tikzcd}
\]
where we agree to take 0 as origin, and to orient from 0 toward $a$), the two functors in question are (unless I am mistaken) the functors suspension $\Sigma$ and ``cosuspension'' or ``loop space'' functor $\Omega$. Subject to checking that $\Omega F$ does indeed coincide with what I have described elsewhere, concerning substantial cartesian squares in $\mathcal{A}$.

In $\Hot_0$ we have a product operation, in fact this must be the categorical product in $\Hot_0$. But on the models it is given by
\[
\mathrm{hot}_0(X) \times \mathrm{hot}_0(Y) \ \simeq \ \mathrm{hot}_0(X \times Y)
\tag*{(8)}
\] % [sic] (8) -> (9), le numéro (8) est déjà employé plus haut sur cette page